\subsection{Rechenregeln}
\label{subsec:vector_algebra_cross_product_rules}

\begin{TheoremBox}{Rechenregeln des Kreuzprodukts}{vector_algebra_cross_product_laws}
  Für $\vec a,\vec b,\vec c\in\R^3$ und $\lambda\in\R$ gilt
  \[
    \begin{aligned}
      \vec a\times\vec b
        &=\Cloze{7cm}{-\vec b\times\vec a}
        &&\text{(Antikommutativgesetz)},\\[0.7em]
      \lambda(\vec a\times\vec b)
        &=\Cloze{10cm}{(\lambda\vec a)\times\vec b
          =\vec a\times(\lambda\vec b)},\\[0.7em]
      \vec a\times(\vec b+\vec c)
        &=\Cloze{9cm}{\vec a\times\vec b+\vec a\times\vec c}
        &&\text{(Distributivgesetz)},\\[0.7em]
      (\vec a+\vec b)\times\vec c
        &=\Cloze{9cm}{\vec a\times\vec c+\vec b\times\vec c}
        &&\text{(Distributivgesetz)}.
    \end{aligned}
  \]
\end{TheoremBox}

Damit begegnen uns erstmals bei einer Vektoroperation zwei vertraute
Rechengesetze, die nicht gelten.

\begin{RemarkBox}[title=Weder kommutativ noch assoziativ]
  Das Kommutativgesetz wird durch die Antikommutativität ersetzt:
  \[
    \vec a\times\vec b=-\vec b\times\vec a.
  \]
  Beide Klammerungen eines dreifachen Kreuzprodukts sind zwar definiert, im
  Allgemeinen gilt jedoch
  \[
    \vec a\times(\vec b\times\vec c)
    \ne(\vec a\times\vec b)\times\vec c.
  \]
\end{RemarkBox}

\begin{ExampleBox}{Reihenfolge beim Kreuzprodukt}{vector_algebra_cross_product_order_example}
  Für
  \[
    \vec a=\begin{pmatrix}1\\2\\3\end{pmatrix},\quad
    \vec b=\begin{pmatrix}4\\5\\6\end{pmatrix},\quad
    \vec c=\begin{pmatrix}7\\8\\9\end{pmatrix}
  \]
  zeigt die vertauschte Reihenfolge
  \[
    \vec a\times\vec b
      =\Cloze{4cm}{\begin{pmatrix}-3\\6\\-3\end{pmatrix}},
    \qquad
    \vec b\times\vec a
      =\Cloze{4cm}{\begin{pmatrix}3\\-6\\3\end{pmatrix}}.
  \]
  Die beiden möglichen Klammerungen liefern
  \[
    \begin{aligned}
      \vec a\times(\vec b\times\vec c)
        &=\Cloze{5cm}{\begin{pmatrix}-24\\-6\\12\end{pmatrix}},\\[0.8em]
      (\vec a\times\vec b)\times\vec c
        &=\Cloze{5cm}{\begin{pmatrix}78\\6\\-66\end{pmatrix}}.
    \end{aligned}
  \]
  Somit ist das Kreuzprodukt weder kommutativ noch assoziativ.
\end{ExampleBox}
